\section{Introduction}
\label{sec:introduction}

A convex body in $\R^n$ is a compact convex set with nonempty interior.
If $K=-K$, then $0$ lies in its interior.
Its support function is
\[
 h_K(u)=\max_{x\in K}\langle x,u\rangle,\qquad u\in S^{n-1}.
\]
For an origin-symmetric body this function is strictly positive.
For a positive continuous function $f$ on the unit sphere, its
\emph{Wulff body} is
\[
 \mathcal W[f]=\bigcap_{u\in S^{n-1}}
 \{x\in\R^n:\langle x,u\rangle\le f(u)\}.
\]
It is a convex body: it contains the ball of radius $\min f>0$ and
is contained in the ball of radius $\max f$. Write $|K|$ for
$n$-dimensional Lebesgue measure. Our main result is the following.

\begin{theorem}[Even logarithmic Brunn--Minkowski inequality]
\label{thm:main}
Let $n\ge1$ and let $K,L\subset\R^n$ be convex bodies satisfying
$K=-K$ and $L=-L$. For every $0\le\lambda\le1$,
\begin{equation}\label{eq:main}
 \left|\mathcal W[h_K^{1-\lambda}h_L^\lambda]\right|
 \ge |K|^{1-\lambda}|L|^\lambda.
\end{equation}
\end{theorem}

The interpolation in \eqref{eq:main} takes the geometric mean of the
support bounds before intersecting the half-spaces. This gives a
strengthening, in the symmetric setting, of the multiplicative form
of the classical Brunn--Minkowski inequality: the arithmetic--geometric
mean inequality gives
\[
 \mathcal W[h_K^{1-\lambda}h_L^\lambda]
 \subset (1-\lambda)K+\lambda L.
\]
In general $h_K^{1-\lambda}h_L^\lambda$ need not itself be a support
function. The intersection defining $\mathcal W$ is essential throughout
the argument.

For $0<p<1$, the $L_p$ interpolation replaces the geometric mean of the
support bounds by the weighted power mean
$((1-\lambda)h_K^p+\lambda h_L^p)^{1/p}$, still taking the Wulff body.
The logarithmic endpoint implies the full symmetric volume inequality in
this intermediate range. B\"or\"oczky, Lutwak, Yang and Zhang record
this monotone implication
\cite[equation~(1.10) and the following paragraph, p.~1977]{BLZY2012}.
The additive
volume-power form below follows by normalizing the bodies and reweighting
the interpolation parameter.

\begin{corollary}[Symmetric $L_p$ Brunn--Minkowski inequality]
\label{cor:lp}
Let $n\ge1$, let $K,L\subset\R^n$ be full-dimensional origin-symmetric
convex bodies, and let $0<p<1$. For every $0\le\lambda\le1$,
\begin{equation}\label{eq:lp}
 \left|\mathcal W\!\left[
  \bigl((1-\lambda)h_K^p+\lambda h_L^p\bigr)^{1/p}
 \right]\right|^{p/n}
 \ge (1-\lambda)|K|^{p/n}+\lambda|L|^{p/n}.
\end{equation}
\end{corollary}

The proof of Corollary~\ref{cor:lp} is given in
Section~\ref{subsec:lp-proof}.


The logarithmic volume theorem has a separate measure-theoretic consequence.
An even log-concave density has the form $e^{-V}$, where $V$ is an even
convex function with values in $\R\cup\{+\infty\}$; no probability
normalization is required. Saroglou proved that the Lebesgue inequality in
dimension $n$ implies the same Wulff inequality for every such density in
that same dimension \cite[Theorem~3.1]{Saroglou2016}.
Taking two scalar dilates of a single body then gives the $(B)$-conjecture:
for the measure $d\mu=e^{-V(x)}\,dx$ and every origin-symmetric convex body
$K$, the function $t\mapsto\mu(e^tK)$ is log-concave on $\R$
\cite[Corollary~3.2]{Saroglou2016}.
Section~\ref{sec:measure} applies this transfer and also treats even
log-concave Radon measures supported on proper subspaces.

\subsection*{Context and prior work}
Theorem~\ref{thm:main} resolves the origin-symmetric logarithmic
Brunn--Minkowski conjecture of B\"or\"oczky, Lutwak, Yang and Zhang
\cite[Problem~1.1]{BLZY2012}.
Stancu \cite[Theorem~2, p.~3264]{Stancu2018} announced the
all-dimensional logarithmic Minkowski inequality and explicitly concluded
the equivalent logarithmic Brunn--Minkowski inequality; the report deferred
details of the flow asymptotics.
The conjecture belongs to the $L_p$ Brunn--Minkowski theory initiated
by Firey's $p$-means of convex bodies
\cite{Firey1962} and developed by Lutwak through $L_p$ mixed volumes
and the $L_p$ Minkowski problem \cite{Lutwak1993}.
For $p>0$, the support bounds are interpolated by the weighted power
mean appearing in Corollary~\ref{cor:lp}; the geometric mean is
its limit as $p\downarrow0$. This endpoint has a particularly close
connection with cone-volume measures and the logarithmic Minkowski
problem. The cone-volume measure records the volumes of cones from the
origin to boundary patches, indexed by their outer normals. B\"or\"oczky, Lutwak, Yang and Zhang established the equivalence
between the corresponding logarithmic volume and mixed-volume inequalities
\cite[Lemma~3.2]{BLZY2012}; their subsequent work characterized which even
measures arise as cone-volume measures \cite[Theorem~1.1]{BLYZ2013}. The latter is an existence result,
separate from the volume inequality proved here.

Several important classes were established in earlier work.
B\"or\"oczky, Lutwak, Yang and Zhang proved the planar case
\cite[Theorem~1.3]{BLZY2012}.
Saroglou established the inequality and studied its equality cases for
bodies unconditional with respect to the same orthonormal basis, meaning that both bodies are invariant under
every coordinate sign change in that basis
\cite[Theorem~1.2]{Saroglou2015}.
The underlying volume inequality for coordinatewise products has antecedents
in Uhrin, Bollob\'as--Leader, and Cordero-Erausquin--Fradelizi--Maurey
\cite{Uhrin1994,BollobasLeader1995,CFM2004}; Saroglou connected this
product to the logarithmic combination. This argument uses the multiplicative
form of the Pr\'ekopa--Leindler inequality.
B\"or\"oczky and Kalantzopoulos extended the symmetry method to bodies
invariant under the same $n$ linear reflections whose fixed hyperplanes
intersect only at the origin \cite[Theorem~2]{BK2022}.
For unit balls of complex norms, Rotem derived the inequality from
Cordero-Erausquin's volume inequality for complex interpolation
\cite{Rotem2014,Cordero2002}.

A complementary approach studies the second variation of volume.
Colesanti, Livshyts and Marsiglietti established a local result near
Euclidean balls and developed infinitesimal formulations
\cite{CLM2017}. Kolesnikov and Milman expressed the local $L_p$ inequality
as a lower bound for the first nonconstant even eigenvalue of the
Hilbert--Brunn--Minkowski operator, and proved it for
$p\ge 1-c n^{-3/2}$ with a universal $c>0$ \cite{KM2022}.
Their spectral formulation explains why symmetry matters: the translation
eigenmodes are odd and hence disappear in the even subspace. Chen, Huang, Li and Liu obtained the corresponding global
inequalities for $p$ sufficiently close to $1$ by a PDE argument
\cite{CHLL2020}. Putterman proved the equivalence of local and global
$L_p$ inequalities for every $p\in[0,1)$ by studying strongly isomorphic
polytopes \cite[Theorem~1.7]{Putterman2021}.
Van Handel proved the local logarithmic inequality for origin-symmetric
zonoids, that is, Hausdorff limits of sums of segments, against arbitrary
origin-symmetric comparison bodies
\cite[Theorem~1.4 of the cited arXiv version]{vanHandel2023}.
Milman's centro-affine interpretation of the operator led to further
curvature-dependent results and an isomorphic logarithmic Minkowski
theorem: every symmetric body has a suitable replacement within a universal
geometric factor \cite{Milman2025}. These developments identify the
analytic obstruction without imposing a common normal fan on the global
conjecture.

Iffland gave another spectral proof of the local logarithmic inequality
for origin-symmetric bodies of revolution and extended it to nonsymmetric
comparison bodies under a weighted centering condition
\cite[Theorem~A]{Iffland2026}.
Lu proved the global inequality for origin-symmetric pairs sharing
$n-2$ orthogonal reflection symmetries
\cite[Theorems~1.2 and~1.3]{Lu2026}.
Xi \cite[Theorem~1.1]{Xi2026} gives a new planar proof that also treats
Dar's conjecture. These results illustrate the different roles of
local variation and additional symmetries in the problem.

\subsection*{Proof strategy}
For finitely many spanning unit vectors $p_i$ and positive numbers $h_i$, the
constraints $|\langle p_i,x\rangle|\le h_i$ define a symmetric
polytope. Replace its indicator by the smooth weight
\[
 \exp\left(-\frac{\eps}{2}|x|^2
       -\sum_i\left(\frac{\langle p_i,x\rangle}{h_i}\right)^q\right),
 \qquad \eps>0,\quad q\ge2\text{ even}.
\]
As $q$ tends to infinity, this weight becomes the Gaussian weight
restricted to the polytope. Letting $\eps$ tend to zero recovers its
volume. Thus it suffices to prove that the integral of the smooth
weight is log-concave when each $h_i$ varies geometrically.
Section~\ref{sec:reduction} makes this reduction precise, including
the passage from finitely many directions to all directions.

Geometric changes of slab widths are also related to the $(B)$-property,
which asks for log-concavity of measure under exponential dilations.
Saroglou showed that the diagonal $(B)$-property for uniform measures on
cubes in every dimension is equivalent to the logarithmic
Brunn--Minkowski conjecture in every dimension
\cite[Theorem~1.5]{Saroglou2015}. This relates the geometric reduction to
an earlier measure-theoretic formulation. The diagonal statement here is
restricted to uniform measures on cubes; the consequence for general even
log-concave measures concerns scalar dilations. The summandwise variance
estimate below supplies the required concavity directly.

The second derivative of the logarithm of that integral is a variance
minus an expected second derivative. The analytic task is therefore a
variance bound for sums of even powers of linear forms
(Proposition~\ref{prop:variance}). We prove it in moment coordinates:
the probability with density proportional to the displayed weight is
the image of $e^{-\varphi(z)}\,dz$ under $x=\nabla\varphi(z)$.
Section~\ref{sec:moment} constructs these coordinates with a bounded,
everywhere positive Hessian, starting from the compact moment theorem
of Berman and Berndtsson \cite{BermanBerndtsson2013} and Klartag's
Hessian estimate \cite{KlartagArxiv2013}. The general moment-measure
theory is due to Cordero-Erausquin and Klartag \cite{CEK2015}.

Two steps complete the analytic argument. First, a differential
operator associated with the moment Hessian produces a dense class
of centered even test functions (Sections~\ref{sec:identities}
and~\ref{sec:density}). The only obstructions to density are affine functions: centering removes
constants, and evenness removes linear functions. Second, a tensor calculation for
these tests has a remainder that is the sum of explicit squares
(Section~\ref{sec:tensor}). This supplies the variance bound by two
Cauchy--Schwarz inequalities (Section~\ref{sec:variance}).

The analytic setting belongs to the tradition of Brascamp--Lieb
variance inequalities for log-concave measures \cite{BrascampLieb1976}.
Kolesnikov, Livshyts and Rotem \cite[Theorem~1.4 and Corollaries~1.5--1.6]{KLR2026}
relate strengthened Brascamp--Lieb inequalities for homogeneous density
potentials to local $p$-Brunn--Minkowski inequalities and obtain the local
logarithmic inequality for $\ell_q$ balls, $q\ge1$.
Their homogeneous variance formulation uses the Hessian of the density
potential; our summandwise estimate uses the separately constructed
moment potential. The transport-Hessian diffusion and curvature formulas
were studied earlier by Kolesnikov \cite{Kolesnikov2014}; the Bochner
identities in \cite[Section~4.3]{KLR2026} build on this framework.

The ingredient here is a tensor estimate in moment coordinates that
controls each homogeneous summand separately. The accompanying density
argument uses a global upper Hessian bound and local strict positivity;
it does not require a global lower Hessian bound. We state these two
steps separately so that their hypotheses and their possible use for
other variance estimates are explicit.
